\section*{Chapter \ref{ch:b2:moral-hazard} --- Moral Hazard and Contracts}

\begin{solution}{exo:b2:moral-hazard:1}
% figdata: safelite_summary
% ledger: b2:lazear2000-t3
The raw means mix months and regions: demand for windscreens varies with the season and the weather, and the firm moved
region by region, so the piece-rate months are not like the hourly months. Controlling for month and year raises the
gain to 44 per cent. The gap between 44 and the 22 per cent found within the same worker, 22 points, is the part due to
who worked at the firm: more productive installers stayed or joined.
\end{solution}

\begin{solution}{exo:b2:moral-hazard:2}
% code: exercises/solutions.py two_outcome
$\sqrt{w} - 2 = 6$ gives $w = 64$ and profit $84 - 64 = 20$. Easy work needs $\sqrt{w} = 6$, $w = 36$, with profit $52 - 36 = 16$.
\end{solution}

\begin{solution}{exo:b2:moral-hazard:3}
% code: exercises/solutions.py hie
$1 - 2.787/3.555 = 21.6$ per cent; $1 - 2.112/2.787 = 24.2$ per cent. Because plans were assigned at random, people in each plan
were alike on average in health and in their taste for care; the differences in use come from the plans.
\end{solution}

\begin{solution}{exo:b2:moral-hazard:4}
% code: exercises/solutions.py team
Each keeps a quarter of what her effort adds, so $e = 0.25$. Her share of output is 0.25 and her cost $0.25^2/2 = 0.03125$: a surplus
of 0.21875, against $1 - 0.5 = 0.5$ at the efficient effort of 1.
\end{solution}

\begin{solution}{exo:b2:moral-hazard:5}
% code: exercises/solutions.py two_outcome
$u_H - u_L = 5$ and $0.8u_H + 0.2u_L = 8$ give $u_L = 4$, $u_H = 9$: $w_L = 16$, $w_H = 81$. The expected wage is $0.8 \times 81 + 0.2 \times 16 = 68$, an
agency cost of $68 - 64 = 4$. Profit from hard work is $84 - 68 = 16$, the same as from easy work: the owner is indifferent.
\end{solution}

\begin{solution}{exo:b2:moral-hazard:6}
% code: exercises/solutions.py linear
$b^* = 1/(1 + 2 \times 1 \times 0.5) = 0.5$ and $e = 0.5$. \omterm{def:b2:expected-utility:risk-premium}{Risk premium} $2 \times 0.25 \times 0.5/2 = 0.125$; cost of effort $0.5^2/2 = 0.125$; surplus
$0.5 - 0.125 - 0.125 = 0.25$; agency cost $0.5 - 0.25 = 0.25$.
\end{solution}

\begin{solution}{exo:b2:moral-hazard:7}
% code: exercises/solutions.py holdup
The supplier invests if its share $s$ of the gain covers the cost: $50s \ge 30$, so $s \ge 0.6$. A fixed price helps only if the
contract can say what is to be delivered; if the design or quality cannot be described or verified, the buyer can refuse
delivery and reopen the bargain once the plant is adapted.
\end{solution}

\begin{solution}{exo:b2:moral-hazard:8}
Plans were assigned at random, so the people in each plan did not choose it and did not differ in health on average: the
differences in use are responses to the price of care. Among ordinary buyers, those who expect to need care choose more
cover, so more use under fuller cover would mix that selection with the response to price.
\end{solution}

\begin{solution}{exo:b2:moral-hazard:9}
% code: exercises/solutions.py rpe_k
$V(0) = 3 + 1 = 4$; $V(1) = 1 + 1 = 2$; $k^* = 3/4 = 0.75$ and $V(0.75) = 3 \times 0.0625 + 1 + 0.5625 = 1.75$. The shares are $1/(1 + 2 \times 4) = 0.111$,
$1/(1 + 2 \times 2) = 0.2$ and $1/(1 + 2 \times 1.75) = 0.222$.
\end{solution}

\begin{solution}{exo:b2:moral-hazard:10}
The first-best surplus is $1/(2c)$ and the second best $b^*/(2c)$ (\cref{prop:b2:moral-hazard:linear}), so the agency cost is $(1 - b^*)/(2c)$.
Since $b^* = 1/(1 + rc\sigma^2)$ falls with $\sigma^2$ and tends to 0, the agency cost rises with $\sigma^2$ and tends to $1/(2c)$.
\end{solution}

\begin{solution}{exo:b2:moral-hazard:11}
% code: exercises/solutions.py risk_neutral
Incentive: $0.4(w_H - w_L) \ge 2$, so $w_H - w_L \ge 5$; participation: $0.8w_H + 0.2w_L - 2 \ge 5$. Any pair with $w_H - w_L \ge 5$ and expected wage 7
works, for instance $w_L = 3$, $w_H = 8$. The first-best wage is also 7 ($w - 2 = 5$), so the agency cost is zero. A risk-neutral agent needs
no insurance, so making her bear the risk costs nothing; in the limit the owner could rent her the orchard for a fixed sum.
\end{solution}

\begin{solution}{exo:b2:moral-hazard:12}
% code: exercises/solutions.py rpe_k
The measure's noise is $(1 - k)\eta + \varepsilon_1 - k\bar\varepsilon$, with variance $V(k) = (1 - k)^2\sigma_\eta^2 + \sigma^2 + k^2\sigma^2/n$. Then $V'(k) = 0$ at $k^* = \sigma_\eta^2/(\sigma_\eta^2 +
\sigma^2/n)$, and $V(k^*) = \sigma^2 + \sigma_\eta^2(\sigma^2/n)/(\sigma_\eta^2 + \sigma^2/n)$. With 3, 1 and 4: $k^* = 3/3.25 = 0.923$ and $V = 1 + 0.75/3.25 = 1.231$. As $n$ grows, $k^* \to 1$
and $V \to \sigma^2$: the average of many others measures the common shock exactly.
\end{solution}

\begin{solution}{pb:b2:moral-hazard:1}
% code: exercises/solutions.py orchard
\textbf{Part I.}
1. $e = b/0.01 = 100b$.
2. $e = 50$, so expected output 50 (hundred dalens) and cost $0.01 \times 2{,}500/2 = 12.5$.
3. Variance $1{,}800 + 200 = 2{,}000$; \omterm{def:b2:expected-utility:risk-premium}{risk premium} $0.05 \times 0.25 \times 2{,}000/2 = 12.5$.
4. $a + be - 0.005e^2 - 50b^2$; at $e = 100b$ this is $a + 100b^2 - 50b^2 - 50b^2 = a$.
5. It is paid whatever the effort.
\textbf{Part II.}
6. $e = 1/0.01 = 100$.
7. $100 - 0.01 \times 10{,}000/2 = 50$.
8. No effort and no surplus.
9. $e = 100$, \omterm{def:b2:expected-utility:risk-premium}{risk premium} $0.05 \times 2{,}000/2 = 50$, surplus $100 - 50 - 50 = 0$.
10. A fixed wage gives no incentive; a fixed rent loads all the risk on a risk-averse tenant.
\textbf{Part III.}
11. $S(b) = 100b - 50b^2 - 50b^2 = 100b - 100b^2$.
12. $b^* = 1/(1 + 0.05 \times 0.01 \times 2{,}000) = 0.5$.
13. Effort and expected output 50.
14. 12.5, 12.5 and $50 - 25 = 25$.
15. $50 - 25 = 25$.
16. $e = 75$; $S = 75 - 100 \times 0.5625 = 18.75$.
17. In the model the share trades incentive against risk and nothing else, and 0.5 is already the best trade. The reform also
gave security of tenure, which encourages effort and investment; and before it, shares may have been held below what
this model would choose, by landowners' threats of eviction among other things.
\textbf{Part IV.}
18. $1{,}800/2{,}000 = 0.9$.
19. Variance $200 + 200 = 400$; $b = 1/(1 + 0.0005 \times 400) = 1/1.2 = 0.833$.
20. $e = 83.3$; surplus $(1/1.2)/0.02 = 41.67$; agency cost $50 - 41.67 = 8.33$.
21. $k^* = 1{,}800/2{,}000 = 0.9$; variance $200 + 1{,}800 \times 200/2{,}000 = 380$.
22. $b = 1/1.19 = 0.840$; $e = 84.0$; surplus $0.8403/0.02 = 42.02$; agency cost 7.98.
23. $S(b) = 100b - b^2(50 + 0.025 \times 380) = 100b - 59.5b^2$; at 0.75: $75 - 33.47 = 41.53$.
24. The tenants could agree to work less together; one could spoil the other's crop; the two orchards may not share the
same weather or soil; and the neighbour's effort may change once the comparison is known.
25. Without comparison the optimal share is one half, effort 50 and the agency cost 25 hundred dalens a season, half the
first-best surplus of 50; comparing with the neighbour at the best weight 0.9 raises the share to 0.84 and cuts the agency
cost to 7.98.
\end{solution}
